\chapter{Drawing the Game: Graphics, Collision, and the Frame}
\label{ch:asm3}

You can now move an object with the joystick and keep it on the screen. This
chapter gives that object a body. We move from the text console to the
320-by-200 graphics viewport, draw rectangles for a paddle and a ball, detect
when they touch, and tie the three functions together into a complete, playable
graphical Pong game. \citep{prg32qemu,prg32tutorialgraphic}.

\section{A blank frame}

We start with a \code{draw} that clears the screen to black and does nothing else: confirm it runs.

\lstinputlisting[
    style=asm,
    caption={%
        \fname{hello\_world\_graphics.S} \protect\\
       A \code{draw} that clears the frame. Build this first and confirm%
       a black screen before drawing anything on it.%
    }
]{examples/ch05/hello_world_graphics.S}

\section{Drawing a rectangle}

The mental model from ASCII based programs is unchanged: state lives in \code{.data},
\code{update} changes it, \code{draw} renders it. What changes is the drawing
call. Instead of \code{prg32\_console\_write}, we use \code{prg32\_gfx\_clear} to
erase the frame and \code{prg32\_gfx\_rect} to draw a filled rectangle.

To draw a rectangle, we use the call \code{prg32\_gfx\_rect(x, y, w, h, color)}. 
In assembly the five arguments go, in order, into \reg{a0} through \reg{a4}, 
following the calling convention. The colour in \reg{a4} is an RGB565 value; 
you may write it as a named constant in C,
but in assembly it is usually a plain number: \code{65535} for white,
\code{65504} for yellow, \code{0} for black \citep{prg32repo}.

\begin{prgwarn}
The five arguments of \code{prg32\_gfx\_rect} occupy \reg{a0}=x, \reg{a1}=y,
\reg{a2}=width, \reg{a3}=height, \reg{a4}=colour. Getting an argument into the
wrong register draws the wrong rectangle: a vivid, immediate way to see what the
calling convention actually means.
\end{prgwarn}

In the next full example we will draw a white paddle near the bottom of the viewport,  
at a position read from \code{.data}. The other dimensions are constant: we will draw
a 64-by-8 white bar at $y=188$. Wide enough to be easy to see in the emulator and
perfect for recreating the old classic Pong. \citep{prg32repo}. Here is what
the paddle draw section will look like:

\lstinputlisting[
    style=asm,
    caption={%
        \fname{paddle\_draw.S} \protect\\
       Drawing the paddle from a stored position. The five arguments fill \reg{a0}--\reg{a4} in order.%
    }
]{examples/ch05/paddle_draw.S}

\section{The ball: position, velocity, and the bounce}

Now we need a ball. The ball needs four words of state: its $x$ and $y$ position
and its $x$ and $y$ velocity (which the examples call \code{dx} and \code{dy}).
Each frame, the velocity is added to the position; when the ball reaches an edge,
the velocity is negated so it travels back the other way. \citep{prg32repo}.

\lstinputlisting[
    style=asm,
    caption={%
        \fname{ball\_motion.S} \protect\\
       Autonomous ball motion on the x and y axis with edge bounce.%
    }
]{examples/ch05/ball_motion.S}

\begin{prgconcept}
\textbf{Velocity} is just a number added to position each frame. A \emph{bounce}
is negating that number when the object reaches a boundary: \code{neg t4, t4}
turns $+2$ into $-2$, so the next frames move the ball back the way it came. With
this one idea: add velocity and negate at the edges, you can make anything bounce.
\end{prgconcept}

\section{Did they touch? Rectangle collision}

A ball that passes straight through the paddle is no fun. We need to know when
the ball's rectangle overlaps the paddle's rectangle. The framework provides
exactly this test: \code{prg32\_sprite\_hitbox(ax, ay, aw, ah, bx, by, bw, bh)}
takes the position and size of two rectangles and returns non-zero when they
overlap \citep{prg32tutorialgraphic}. In assembly the eight arguments fill
\reg{a0} through \reg{a7}.

\lstinputlisting[
    style=asm,
    caption={%
        \fname{collision\_detection.S} \protect\\
       Testing ball-versus-paddle overlap. A non-zero result in \reg{a0} means they touched this frame.%
    }
]{examples/ch05/collision_detection.S}

Two rectangles overlap when they intersect on \emph{both} axes at once: the $x$
ranges overlap \emph{and} the $y$ ranges overlap. \code{prg32\_sprite\_hitbox}
performs this two-axis test for you and returns the answer as a single value, so
collision becomes one call followed by one branch.

\begin{prgconcept}
Notice that the bounce reuses the clamp from Chapter~\ref{ch:asm2}, with one
addition: where the clamp merely stopped the object at the edge, the bounce also
negates the velocity. The platform's examples keep the clamp so the ball never
visibly leaves the screen even for a single frame, then negate so it returns. Good
game feel often comes from such small combinations of simple rules.
\end{prgconcept}

\section{Assembling the complete game}

We now have every part: a paddle moved by input and clamped, a ball that moves
and bounces, and a collision test. Putting them in the three functions gives a
complete graphical game whose structure exactly matches the platform's
\code{pong\_graphics} example \citep{prg32repo}. The \code{update} moves the 
paddle from input,  advances and bounces the ball, and applies collision; 
the \code{draw} clears the frame and renders the paddle and ball.

\lstinputlisting[
    style=asm,
    caption={%
        \fname{pong\_skeleton.S} \protect\\
       The \code{draw} of the complete game. Compare it line for line%
        with \code{pong\_graphics\_draw} in the repository.%
    }
]{examples/ch05/pong_skeleton.S}

A brief stylistic note: earlier chapters tended to use simple numeric 
labels for branch targets; in this chapter we switch to descriptive, 
named labels for different sections of control flow. This is a readability 
improvement only; it does not change program behaviour but makes the 
structure and intent of each labelled section easier to follow.

\begin{prgtargets}
\textbf{Both targets.} Develop and watch this on the QEMU virtual screen, 
where the ball bounces and the paddle responds to injected input. Then 
build the identical source for the ESP32\nobreakdash-C6 to play it with 
a real joystick and hear the beep. Because QEMU's defaults disable the 
physical GPIO buttons, the paddle may sit still in the emulator unless 
input is injected; the ball still moves, so the render path is fully 
verified on the desktop \citep{prg32qemu}.
\end{prgtargets}

\section{Step-by-step: how this game was built}

It is worth making explicit the pedagogical order in which a game like this is
constructed, because it is the order the platform recommends and the order in
which the example games in \code{examples/games} were themselves written
\citep{prg32tutorial,prg32tutorialgraphic}.

First, write \code{init}, \code{update}, and \code{draw} that only return, 
and confirm the framework runs them. Second, add one \code{.data} variable.
Third, read input in \code{update} and move the objects, then add clamping. 
Fourth, add the collision test. Only then, fifth, add sound on the collision 
event. Each step is buildable and observable on its own; if a step breaks, 
the bug is in the few lines you just added, not hidden in a wall of untested code.

\begin{prgtry}
Build the game incrementally in exactly these five steps, rebuilding and running
after each. Keep a short log: after each step, write one line saying what you
added and what you saw. This log is not busywork: it is the written form of the
``change one thing, observe, record'' method, and it will make the difference
between a frustrating afternoon and a productive one \citep{prg32repo}.
\end{prgtry}

\begin{prgcheck}
You can now clear the frame, draw rectangles from stored state, move the ball
with velocity and bounces, and detect collisions with one call. You have, in
other words, built a complete graphical game in \riscv assembly that runs on both
targets. Chapter~\ref{ch:asmcap} turns this into a capstone you extend and
package as a cartridge.
\end{prgcheck}
